\documentclass[12pt]{article}
\usepackage[utf8]{inputenc}
\usepackage[vietnamese]{babel}
\usepackage{geometry}
\usepackage{fancyhdr}
\usepackage{amsmath}

% Cấu hình Header và Footer
\pagestyle{fancy}
\fancyhf{}
\fancyhead[L]{Xử Lý tín hiệu số}
\fancyhead[C]{HK261}
\fancyhead[R]{CO2036}
\fancyfoot[R]{\thepage}
\renewcommand{\headrulewidth}{0pt}
\renewcommand{\footrulewidth}{0pt}
\setlength{\headheight}{15pt}

% Cấu hình lề trang (chuẩn A4, lề 2cm)
\geometry{a4paper, margin=2cm}

% Định dạng đánh số section có dấu chấm (1., 2., ...) và subsection (4.1, 4.2, ...)
\makeatletter
\def\@seccntformat#1{\csname the#1\endcsname\ }
\makeatother
\renewcommand{\thesection}{\arabic{section}.}
\renewcommand{\thesubsection}{\arabic{section}.\arabic{subsection}}

% Các gói bổ trợ thường dùng cho bài thực hành Xử Lý Tín Hiệu Số (DSP)
\usepackage{amssymb}           % Công thức toán học (Z-transform, DFT, FFT,...)
\usepackage{graphicx}         % Chèn hình ảnh đồ thị tín hiệu
\usepackage{float}            % Cố định vị trí hình ảnh với [H]
\usepackage{listings}         % Chèn mã nguồn SciLab / MATLAB / Python
\usepackage{xcolor}           % Màu sắc
\usepackage{hyperref}         % Chèn liên kết
\hypersetup{colorlinks=true, linkcolor=black, urlcolor=blue}

% Style mã nguồn Scilab (dùng: \begin{lstlisting} ... \end{lstlisting})
\lstdefinestyle{scilab}{
    language=Matlab,
    morecomment=[l]{//},
    basicstyle=\ttfamily\small,
    keywordstyle=\color{blue},
    commentstyle=\color{gray},
    stringstyle=\color{red!60!black},
    numbers=left, numberstyle=\tiny, numbersep=6pt,
    frame=single, breaklines=true, columns=fullflexible
}
\lstset{style=scilab}

% Thông tin báo cáo
\title{BÀI THỰC HÀNH XỬ LÝ TÍN HIỆU SỐ -- LAB 4}
\author{Nguyễn Minh Đức - 2410816 \\ Trần Gia Huân - 2533019}
\date{24/09/2026}

\begin{document}
\maketitle
\thispagestyle{fancy}

\section*{Tổng Quan}

\section{Giới Thiệu}
Lab 4 phân tích tín hiệu và hệ LTI rời rạc trong miền $z$, đồng thời hiện thực tích chập tuyến tính / tích chập vòng trên Scilab. Báo cáo gồm hai phần:

\begin{itemize}
    \item \textbf{Biến đổi Z (Exercise 1--3):} tìm $X(z)$ và ROC cho ba lớp tín hiệu: chiều dài hữu hạn, nhân quả (right-sided), hai phía (two-sided).
    \item \textbf{Scilab bổ sung (Exercise 4--7):}
    \begin{itemize}
        \item Ex 4, 5: khảo sát thư viện/toolbox Scilab cho xử lý audio và xử lý ảnh, kèm demo trên ảnh số.
        \item Ex 6: tích chập tuyến tính $y(n) = \sum_k h(k)x(n-k)$ -- điều kiện hội tụ, tính bằng phương pháp gập--dịch và phương pháp ma trận.
        \item Ex 7: tích chập vòng $y(n) = \sum_{k=0}^{M-1} h(k)\,x[(n-k) \bmod N]$ -- cơ sở công thức, tính bằng hai phương pháp tương tự.
    \end{itemize}
\end{itemize}

\noindent Công cụ: Scilab (core + toolbox cài qua ATOMS), \LaTeX{} cho báo cáo.

\section{Mục Tiêu}
\begin{itemize}
    \item Tính $X(z)$ và xác định ROC cho tín hiệu hữu hạn, nhân quả, phản nhân quả và tổ hợp tuyến tính của chúng.
    \item Liên hệ vị trí cực -- ROC -- tính nhân quả / ổn định của tín hiệu.
    \item Nắm các hàm/toolbox Scilab cho audio (đọc, ghi, phát) và ảnh (histogram, cân bằng histogram, lọc làm mờ, watermark).
    \item Hiện thực tích chập tuyến tính và tích chập vòng bằng hai cách: gập--dịch (folding \& shifting) và ma trận; đối chiếu với hàm có sẵn.
    \item Vẽ $x(n)$, $y(n)$ và năng lượng tín hiệu trên cùng một figure; so sánh kết quả tích chập tuyến tính và tích chập vòng.
\end{itemize}

\section{Phương Pháp}

\subsection{Biến đổi Z và ROC}
Định nghĩa (biến đổi Z hai phía):
\[
X(z) = \sum_{n=-\infty}^{\infty} x(n)\,z^{-n}, \qquad \text{ROC} = \{ z \in \mathbb{C} : \textstyle\sum |x(n)z^{-n}| < \infty \}
\]
Các cặp biến đổi dùng trong bài:
\begin{center}
\begin{tabular}{lll}
\hline
$x(n)$ & $X(z)$ & ROC \\
\hline
$\delta(n-k)$ & $z^{-k}$ & $\mathbb{C}$, bỏ $z=0$ nếu $k>0$, bỏ $z=\infty$ nếu $k<0$ \\
$a^n u(n)$ & $\dfrac{1}{1-az^{-1}}$ & $|z| > |a|$ \\
$-a^n u(-n-1)$ & $\dfrac{1}{1-az^{-1}}$ & $|z| < |a|$ \\
\hline
\end{tabular}
\end{center}
Quy trình:
\begin{enumerate}
    \item Tách $x(n)$ thành tổng các thành phần cơ bản, áp dụng tính tuyến tính.
    \item ROC của tổng $\supseteq$ giao ROC các thành phần (bằng nhau nếu không có triệt tiêu cực--không).
    \item Quy đồng để xác định cực/không. ROC không chứa cực; right-sided $\Rightarrow$ ROC ngoài cực lớn nhất, left-sided $\Rightarrow$ ROC trong cực nhỏ nhất, two-sided $\Rightarrow$ vành khuyên (rỗng thì $X(z)$ không tồn tại).
\end{enumerate}

\subsection{Tích chập tuyến tính}
\[
y(n) = \sum_{k=-\infty}^{\infty} h(k)\,x(n-k)
\]
Với $x(n)$ dài $N$, $h(n)$ dài $M$ (cùng bắt đầu tại $n=0$): $y(n)$ dài $N+M-1$, $n \in [0, N+M-2]$.

\noindent\textbf{Gập--dịch (folding \& shifting):} với mỗi $n$: gập $x(k) \to x(-k)$, dịch $n$ mẫu $\to x(n-k)$, nhân từng mẫu với $h(k)$, cộng dồn.

\noindent\textbf{Ma trận:} $\mathbf{y} = \mathbf{X}\mathbf{h}$, với $\mathbf{X}$ là ma trận Toeplitz kích thước $(N+M-1) \times M$, cột $j$ là $x(n)$ dịch xuống $j$ mẫu, phần thiếu điền 0:
\[
\mathbf{X}_{ij} = x(i-j), \quad 0 \le i \le N+M-2,\ 0 \le j \le M-1
\]

\subsection{Tích chập vòng}
\[
y(n) = \sum_{k=0}^{M-1} h(k)\,x[(n-k) \bmod N], \quad 0 \le n \le N-1
\]
$x(n)$ được coi là một chu kỳ của tín hiệu tuần hoàn $\tilde{x}(n) = x(n \bmod N)$, nên chỉ số $n-k$ được quấn về $[0, N-1]$.

\noindent\textbf{Gập--dịch:} gập vòng $x[(-k) \bmod N]$, dịch vòng $n$ mẫu, nhân với $h(k)$, cộng dồn.

\noindent\textbf{Ma trận:} $\mathbf{y} = \mathbf{C}\mathbf{h}$, $\mathbf{C}$ là ma trận vòng (circulant) $N \times M$:
\[
\mathbf{C}_{ij} = x[(i-j) \bmod N]
\]
Quan hệ với tích chập tuyến tính: $y_{\text{vòng}}(n) = \sum_{r} y_{\text{tt}}(n + rN)$, $0 \le n \le N-1$ (chồng lấn miền thời gian). Hai kết quả trùng nhau khi zero-pad lên $L \ge N+M-1$.

\subsection{Năng lượng tín hiệu}
\[
E_x = \sum_{n} |x(n)|^2
\]
Vẽ năng lượng theo từng mẫu $|x(n)|^2$, $|y(n)|^2$ (stem) cùng figure với $x(n)$, $y(n)$; ghi tổng $E_x$, $E_y$ trên tiêu đề subplot.

\subsection{Công cụ Scilab}
\begin{itemize}
    \item Tín hiệu / tích chập: \texttt{conv}, \texttt{convol} (đối chiếu kết quả), \texttt{toeplitz}, \texttt{modulo}, \texttt{sum}.
    \item Đồ thị: \texttt{subplot}, \texttt{plot2d3} (stem), \texttt{xtitle}, \texttt{xs2png} (xuất ảnh).
    \item Audio (core): \texttt{wavread}, \texttt{wavwrite}, \texttt{auread}, \texttt{auwrite}, \texttt{sound}, \texttt{playsnd}.
    \item Ảnh: toolbox cài qua ATOMS (\texttt{atomsInstall("IPCV")}, SIVP).
\end{itemize}

\section{Nội Dung Thực Hành}
\subsection{Exercise 1}
Tìm $X(z)$ và ROC của:
\[
x(n) = 2\delta(n+2) - \delta(n+1) + 2\delta(n) - 3\delta(n-1) + 4\delta(n-2)
\]

% TODO: lời giải
\[
X(z) = 2z^2 - z + 2 - 3z^{-1} + 4z^{-2} = \frac{2z^4 - z^3 + 2z^2 - 3z + 4}{z^2}
\]
\[
ROC: 0 < |z| < \infty
\]
\subsection{Exercise 2}
Tìm $X(z)$ và ROC của:
\[
x(n) = 0.5^n u(n) + 0.4^n u(n)
\]

% TODO: lời giải
\[
X(z) = \frac{1}{1-0.5z^{-1}} + \frac{1}{1-0.4z^{-1}} 
= \frac{2 - 0.9z^{-1}}{1 - 0.9z^{-1} + 0.2z^{-2}}
\]
ROC là giao của $ROC\left(\frac{1}{1-0.5z^{-1}}\right)$ và $ROC\left(\frac{1}{1-0.4z^{-1}} \right)$
\[
ROC = \{|z|>0.5\} \cap \{|z|>0.4\} = \{|z|>0.5\}
\]
\subsection{Exercise 3}
Tìm $X(z)$ và ROC của:
\[
x(n) = 0.5^n u(n) + 0.9^n u(-n-1)
\]

% TODO: lời giải
Ta tách $x(n)$ thành $f(n)$ và $g(n)$, $x(n) = f(n) - g(n)$, ta có:
\begin{center}
    $\begin{cases}
        f(n) = 0.5^nu(n)\\
        g(n) = 0.9^nu(-n-1)
    \end{cases}$
\end{center}
Biến đổi $f(n)$ sang miền Z, ta có:\\
\[
F(z) = \frac{1}{1-0.5z^{-1}}
\]
\[
ROC_F = \{|z| > 0.5\}
\]
Biến đổi $g(n)$ sang miền Z, ta có:\\
\[
g(n) = -0.9^nu(-n-1)
G(z) = \frac{1}{1-0.9z^{-1}}
\]
\[
ROC_G = \{|z| < 0.9\}
\]
Cuối cùng, ta có kết quả của $X(z)$:
\[
X(z) = F(z) - G(z) = \frac{1}{1-0.5z^{-1}} - \frac{1}{1-0.9z^{-1}} = \frac{-0.4z^{-1}}{1 - 1.4^{-1} + 0.45z^{-2}}
\]
\[
ROC = ROC_F \cap ROC_G = {0.5 < |z| < 0.9}
\]

\section{Bàn Luận}
\subsection{Exercise 4}
Khảo sát thư viện / toolbox Scilab cho xử lý dữ liệu audio.

% TODO: lời giải

\subsection{Exercise 5}
Khảo sát thư viện / toolbox Scilab cho xử lý ảnh; demo một thao tác đơn giản trên ảnh số (hiển thị histogram, cân bằng histogram, làm mờ, watermark).

% TODO: lời giải

% \begin{lstlisting}
% // Scilab code
% \end{lstlisting}

% \begin{figure}[H]
%     \centering
%     \includegraphics[width=0.75\textwidth]{img/chap5/ex5.png}
%     \caption{...}
% \end{figure}

\subsection{Exercise 6}
Tích chập tuyến tính: $y(n) = \sum_{k=-\infty}^{\infty} h(k)\,x(n-k)$.
\begin{enumerate}
    \item Rút gọn $y(n)$ để tổng hội tụ trong hai trường hợp:
    \begin{itemize}
        \item $h(n)$ và $x(n)$ cùng nhân quả.
        \item Hệ FIR: $h(n) = 0$ với $n < 0 \lor n > M$.
    \end{itemize}
    \item (Trường hợp 2) Tính $y(n)$ với $x(n) = \{\underset{\uparrow}{1}, 2, -3, 2, 1\}$, $h(n) = \{\underset{\uparrow}{1}, 0, -1\}$ bằng phương pháp gập--dịch. Hiện thực bằng Scilab.
    \item (Trường hợp 2) Tính lại $y(n)$ bằng phương pháp ma trận. Hiện thực bằng Scilab.
    \item Vẽ $x(n)$, $y(n)$ và năng lượng tín hiệu trên cùng một figure.
\end{enumerate}

% TODO: lời giải

% \begin{lstlisting}
% // Scilab code
% \end{lstlisting}

% \begin{figure}[H]
%     \centering
%     \includegraphics[width=0.75\textwidth]{img/chap5/ex6.png}
%     \caption{...}
% \end{figure}

\subsection{Exercise 7}
Tích chập vòng giữa $x(n)$ tuần hoàn chiều dài $N$ và $h(n)$ chiều dài $M$:
\[
y(n) = \sum_{k=0}^{M-1} h(k)\,x[(n-k) \bmod N]
\]
Giải thích cơ sở của công thức trên.
\begin{enumerate}
    \item Tính $y(n)$ với $x(n) = \{\underset{\uparrow}{1}, 2, -3, 2, 1\}$, $h(n) = \{\underset{\uparrow}{1}, 0, -1, -1, 1\}$ bằng phương pháp gập--dịch. Hiện thực bằng Scilab.
    \item Tính lại $y(n)$ bằng phương pháp ma trận. Hiện thực bằng Scilab.
    \item Vẽ $x(n)$, $y(n)$ và năng lượng tín hiệu trên cùng một figure.
\end{enumerate}

% TODO: lời giải

% \begin{lstlisting}
% // Scilab code
% \end{lstlisting}

% \begin{figure}[H]
%     \centering
%     \includegraphics[width=0.75\textwidth]{img/chap5/ex7.png}
%     \caption{...}
% \end{figure}


\section{Tài Liệu Tham Khảo}
\begin{enumerate}
    \renewcommand{\labelenumi}{[\theenumi]}
    \item J. G. Proakis, D. G. Manolakis, \textit{Digital Signal Processing: Principles, Algorithms, and Applications}, 4th ed., Pearson, 2007.
    \item A. V. Oppenheim, R. W. Schafer, \textit{Discrete-Time Signal Processing}, 3rd ed., Pearson, 2010.
    \item Đề bài \textit{Lab \#4 -- Digital Signal Processing}, Trường Đại học Bách Khoa -- ĐHQG TP.HCM.
    \item Scilab Online Help. \url{https://help.scilab.org}
    \item ATOMS -- Image Processing and Computer Vision Toolbox (IPCV). \url{https://atoms.scilab.org/toolboxes/IPCV}
    \item ATOMS -- Scilab Image and Video Processing Toolbox (SIVP). \url{https://atoms.scilab.org/toolboxes/SIVP}
\end{enumerate}


\end{document}

